\section{Introduction}
\label{sec:intro}

Autoregressive language-model decoding reads each active request's key--value (KV) cache on every decode step.
The cache footprint grows with batch size and context length; storing it consumes high-bandwidth memory (HBM) capacity, and streaming it from HBM often dominates decode-attention latency.
Low-bit KV quantization reduces that traffic directly.
KIVI~\cite{liu2024kivi}, Kitty~\cite{xia2025kitty}, OSCAR~\cite{zhou2026oscar}, KVQuant~\cite{hooper2024kvquant}, and QServe~\cite{lin2025qserve} store each cached value in 2--4 bits with a group-shared scale and offset.

Compressing the cache in memory does not eliminate on-chip operand preparation.
KIVI~\cite{liu2024kivi}, QServe~\cite{lin2025qserve}, TurboAttention~\cite{kang2025turboattention}, BitDecoding~\cite{du2026bitdecoding}, SAW-INT4~\cite{jia2026saw}, and OSCAR~\cite{zhou2026oscar} all convert or reconstruct low-bit values before each matrix product.
This preparation consumes instruction-issue throughput, registers, and shared-memory bandwidth even when fused into attention without writing dequantized K and V to HBM.
\emph{Although bit-packing reduces HBM traffic, per-element operand preparation can still limit the resulting latency speedup.}

\sys{} applies affine folding to the query--key (QK) and probability--value (PV) matrix products, which reduce over channels and tokens.
For token-wise K, the kernel scales each group's partial QK result and adds the group's offset times the matching query sum.
For token-wise V, the kernel scales the attention weights before PV and adds a weighted sum of offsets to the output.
Channel-wise K instead scales the query before QK, while channel-wise V scales each group's partial PV result, and both add the corresponding offset-weighted sum.
These per-group operations replace element-wise scale multiplication and offset addition, while bit manipulation constructs Tensor Core KV operands directly from packed integers.
In exact arithmetic, affine folding produces the same attention output as dequantizing every KV value before the matrix products.

This paper makes four contributions:
\begin{itemize}
  \item \textbf{Identify decode attention performance bottleneck.}
        We compare per-byte reconstruction instruction demand against peak instruction-issue and HBM throughputs on H100 and B200, showing how per-element conversion and affine arithmetic limit low-bit decode speedups (\Cref{sec:analysis}).
   \item \textbf{Efficient Tensor Core operand preparation via affine folding.}
        We formulate affine folding for independent token-wise or channel-wise K and V quantization, paired with bit-level Tensor Core operand construction, to feed Tensor Cores efficiently with lower instruction demand (\Cref{sec:design-fold,sec:design-unpack}).
  \item \textbf{Fast GPU kernels for Hopper and Blackwell.}
        We design warp-specialized decode pipelines for H100 register accumulators and B200 tensor-memory (TMEM) accumulators (\Cref{sec:design-pipelines}).
  \item \textbf{Kernel and serving-engine evaluation.}
        We evaluate \sys{} in standalone kernel benchmarks and integrated into SGLang~\cite{zheng2024sglang}.
        \sys{} INT2 achieves geometric-mean kernel-latency speedups over FlashInfer BF16 of $3.36\times$ on H100 and $2.67\times$ on B200.
        For long-context serving workloads, \sys{} INT2 increases peak aggregate serving throughput by up to $170\%$ and $341\%$ over 2-bit and 16-bit baselines, respectively (\Cref{sec:eval}).
\end{itemize}
